% rob_t7_fxgcf_construction — the four FXGCF constructions compared.
% Numbers from the construction-comparison analysis (fxgcf_comparison/,
% build_comparison.R), run on the identical thesis pipeline.
\small
\begin{tabular}{ll c c c c c}
  \toprule
  Construction & Weights & $\Corr(\cdot,\GCF)$ & Mean IS $R^{2}$
    & $|t|>1.96$ & $R^{2}_{\mathrm{oos}}$ (pooled) & OOS$+$ \\
  \midrule
  Bottom-up (baseline) & GDP   & $0.81$ & $0.143$ & $10/11$ & $+0.021$ & $6/11$ \\
  Bottom-up            & $1/n$ & $0.76$ & $0.140$ & $10/11$ & $+0.016$ & $6/11$ \\
  Top-down             & GDP   & $0.99$ & $0.101$ & $5/11$  & $+0.010$ & $9/11$ \\
  Top-down             & $1/n$ & $0.95$ & $0.109$ & $6/11$  & $+0.033$ & $8/11$ \\
  \bottomrule
\end{tabular}
\tabnotes{All four constructions use the same cycle predictors, the one-year
  cycle $\cyc^{(1)}_{i,t}$ and the average cycle $\cycbar_{i,t}$, and predict
  the US-dollar excess return $\rxbar^{\,\mathrm{USD}}_{i,t+12}$. The
  bottom-up variants fit a USD cycle factor per country
  \eqref{eq:fxcf-reg}--\eqref{eq:fxcf} and aggregate the fitted factors; the
  top-down variants first aggregate the dollar return and the predictors and
  fit one regression on the aggregates. Weights are nominal-GDP shares or
  equal ($1/n$). ``Mean IS $R^{2}$'' is the cross-country mean in-sample fit,
  ``$|t|>1.96$'' the number of markets with a significant HAC $t$-statistic,
  and the two out-of-sample columns report the pooled fully-recursive
  $R^{2}_{\mathrm{oos}}$ statistic and the number of markets with a
  positive value.}
